\section{Methods}
\label{sec:methods}

The methodology follows a single pipeline from participant selection to statistical analysis. A primary cohort of 871 participants is first selected from ABIDE-I to reduce the influence of poor-quality scans and provide a more consistent sample for the primary analyses. Selection is based on the available anatomical and functional image-quality ratings from three raters, with participants retained only when the required ratings indicate acceptable, or where permitted borderline, image quality. Preprocessed regional time series are then obtained for each participant, representing changes in the signal of each brain region over the course of the scan.

Three functional-connectivity representations are constructed from these time series. The first captures whole-scan correlations between pairs of brain regions using shrinkage-based covariance estimation. The second represents whole-scan connectivity in tangent space by describing each participant's covariance pattern relative to a reference covariance estimated from the training data. The third extends this approach to shorter windows of the scan, allowing connectivity to be represented across multiple temporal segments. Each representation is standardised and classified separately using $\ell_2$-regularised logistic regression, after which the three predicted probabilities are averaged to obtain the final ensemble probability.

Model performance is evaluated using both pooled ten-fold cross-validation and leave-one-site-out validation, while classifier regularisation is selected through nested cross-validation. All quantities required for model fitting are estimated from the corresponding training data to prevent information from the held-out data from entering the primary pipeline. The resulting out-of-fold predictions are subsequently used to assess discrimination, calibration, and selective prediction. Additional analyses examine the effects of alternative connectivity representations, classifiers, atlases, feature-selection strategies, cohort definitions, and participant subgroups, together with associations between classifier scores and symptom measures. Separate group-contrast analyses evaluate diagnostic-group differences in functional connectivity after adjustment for relevant nuisance variables. Fold-level results and out-of-fold predictions are retained to support paired statistical comparisons across experimental conditions. The complete methodological workflow is summarised in Figure~\ref{fig:overview}.

Figure~\ref{fig:overview} summarises the primary modelling pipeline. A primary cohort of 871 participants is obtained from ABIDE-I after image-quality screening based on the available ratings from three raters. Preprocessed regional time series are then extracted from the C-PAC \texttt{filt\_noglobal} derivative using the 200-region CC200 atlas. From these signals, three complementary functional-connectivity representations are constructed: shrinkage-based whole-scan correlation, static tangent-space connectivity, and windowed tangent-space connectivity computed from six temporal windows of 50 volumes. For the tangent-space representations, reference covariance matrices are estimated exclusively from the corresponding training data to prevent information from the held-out partition from entering feature construction.

Each connectivity matrix is vectorised into 19,900 pairwise features and standardised using training-partition statistics. Separate $\ell_2$-regularised logistic-regression models are then fitted to the three representations, with the inverse regularisation parameter $C$ selected by inner three-fold cross-validation. The resulting probabilities are averaged with equal weights to produce the final ensemble prediction. Generalisation is evaluated using both pooled ten-fold cross-validation and leave-one-site-out validation across the 20 site labels. Throughout the primary pipeline, all data-dependent transformations and model parameters are estimated within the outer training partition before being applied to held-out participants, thereby maintaining separation between model fitting and evaluation.

\begin{figure}[pos=htbp]
\centering
\includegraphics[width=\linewidth,height=0.84\textheight,keepaspectratio]{fig_overview}
\caption{Overview of the primary analysis pipeline. (a) ABIDE-I participants undergo image-quality control to define the primary cohort, followed by extraction of CC200 regional time series from the preprocessed fMRI data. (b) Three functional-connectivity representations are derived for each participant: shrinkage-based correlation, static tangent-space connectivity, and windowed tangent-space connectivity. Tangent-space reference matrices are estimated from the corresponding training partition only. (c) Each representation is vectorised and standardised using training-derived statistics before fitting a separate $\ell_2$-regularised logistic-regression model. The regularisation parameter is selected by inner three-fold cross-validation, and the three model probabilities are averaged to obtain the final ensemble prediction. Generalisation is evaluated using pooled ten-fold and leave-one-site-out cross-validation. Circled numbers indicate the principal stages of the pipeline.}
\label{fig:overview}
\end{figure}

The Methods are organised into nine subsections that follow the analysis pipeline. Section~\ref{subsec:study-design-cohort} describes the study design, data, cohorts, phenotypic variables, atlas, and scan-quality characteristics. Sections~\ref{subsec:connectivity-representations}--\ref{subsec:probability-assessment} then define the connectivity representations, prediction model, validation strategy, and probability-based analyses. The planned comparisons and exploratory group-contrast analyses are presented in Sections~\ref{subsec:planned-comparisons} and~\ref{subsec:group-differences}, followed by the statistical analysis in Section~\ref{subsec:statistical-analysis}. Reproducibility and the scope of the resulting estimates are addressed in Section~\ref{subsec:reproducibility-scope}.


\subsection{Study design and cohort}
\label{subsec:study-design-cohort}

This subsection defines the study question, the data, and the participant cohorts on which
the analyses are based. These definitions are given first, because all representations,
models, and comparisons are based on these data and cohorts. The study design is
described in Section~\ref{subsubsec:study-design}, and the data source and preprocessing in
Section~\ref{subsubsec:data-preprocessing}. The primary cohort is then defined in
Section~\ref{subsubsec:cohorts}, and the wider quality-control cohort in Section~\ref{subsubsec:cohorts}.
Next, the phenotypic variables and the handling of missing values are specified in
Section~\ref{subsubsec:phenotypic-variables}. The atlas and the signal-quality procedure are described in
Section~\ref{subsubsec:atlas-signal-quality}. Finally, the treatment of differences in scan length is described
in Section~\ref{subsubsec:scan-length}.

\subsubsection{Study design}
\label{subsubsec:study-design}

Clinical diagnosis of autism spectrum disorder (ASD) combines direct observation,
developmental history, and professional judgement. In this study, ASD status is defined by
the diagnostic labels provided in the Autism Brain Imaging Data Exchange I (ABIDE-I), while
resting-state functional magnetic resonance imaging (rs-fMRI) is used to characterise
large-scale functional connectivity in autistic and non-autistic participants.

The ASD classification target is derived directly from the ABIDE-I diagnostic labels.
Participants with \texttt{DX\_GROUP = 1} are coded as autistic, whereas those with
\texttt{DX\_GROUP = 2} are coded as non-autistic. The primary analysis therefore evaluates
how well the imaging-based models distinguish autistic from non-autistic participants within
the ABIDE-I research cohort.

Discrimination is evaluated primarily using participant-pooled out-of-fold area under the
receiver operating characteristic curve (AUROC) under pooled ten-fold and
leave-one-site-out validation. For each participant, the prediction used for evaluation is
generated by a model fitted without that participant's held-out data, and predictions from
all held-out folds are subsequently combined to calculate the overall AUROC. AUROC is used
as the primary discrimination measure because it assesses the ranking of autistic and
non-autistic participants across prediction scores without requiring a fixed classification
threshold.

In pooled ten-fold validation, participants from all sites are distributed across ten folds,
with each fold held out once, providing an estimate of generalisation within the combined
ABIDE-I multisite cohort. Leave-one-site-out validation instead holds out all participants
associated with one site label at a time, providing a more stringent assessment of
generalisation to an ABIDE-I site that is absent from model training.

The primary prediction model is an equal-weight ensemble of three connectivity representations. Beyond this model, six comparative analyses address the main contributions of the study. Four of these examine how discrimination varies with cognitive ability, whether functional-connectivity differences persist after matching for head motion, how feature-selection scope and granularity influence predictive performance, and how different classifier families perform under a common representation and validation procedure. The remaining two concern the symptom measures and the windowed description. For the symptom measures, associations with classifier scores are examined both across the full cohort and within autistic participants. For the windowed description, the whole-scan description is given the same covariance shrinkage, so that the two are compared under matched regularisation. Figure~\ref{fig:designs}b--e summarises the designs of the first four comparative analyses, including the corresponding matching, feature-selection, subgroup, and classifier-comparison procedures.


\subsubsection{Data source and preprocessing}
\label{subsubsec:data-preprocessing}

Data are obtained from the Autism Brain Imaging Data Exchange I (ABIDE-I), a multi-site research dataset that combines resting-state functional MRI (rs-fMRI), anatomical MRI, and phenotypic information from 1,112 participants across 17 acquisition sites \citep{DiMartinoetal2014}. The phenotypic records provide diagnostic status together with demographic and cognitive measures, including age, sex, full-scale IQ, verbal IQ, performance IQ, handedness, and acquisition site. These variables are retained to characterise the study population and to support the subgroup and covariate analyses described later. After the quality-control procedure described in Section~\ref{subsubsec:cohorts}, the primary cohort contains 871 participants, including 403 autistic and 468 non-autistic participants. The cohort includes 727 male and 144 female participants and spans 6.5--58 years of age, with a mean age of 16.9 years (SD 7.6). Mean full-scale IQ is 105.8 in the autistic group and 111.1 in the non-autistic group. Because ABIDE-I combines data acquired independently at multiple centres, heterogeneity remains across sites in scanner hardware, acquisition protocols, scan duration, participant composition, and image quality \citep{Abrahametal2017}. Image quality is therefore considered explicitly during cohort selection using the anatomical and functional quality ratings provided with the release. Preprocessed imaging outputs are obtained from the Preprocessed Connectomes Project \citep{Craddocketal2013}, which provides ABIDE-I data processed through standardised workflows and thereby supports consistent preprocessing across participants from different acquisition sites. The use of these publicly released derivatives also improves reproducibility by fixing the imaging inputs and preprocessing choices. The present study uses outputs generated with the Configurable Pipeline for the Analysis of Connectomes (C-PAC), specifically the \texttt{filt\_noglobal} processing strategy.

Within the selected C-PAC workflow, temporal and spatial sources of unwanted variation are addressed before functional connectivity is estimated. Slice-acquisition timing differences are corrected, successive functional volumes are realigned to reduce motion-related displacement, non-brain tissue is removed, signal intensities are normalised, and the functional images are aligned with the corresponding anatomical images. Additional variation associated with head motion and non-neural signals is reduced through nuisance regression using a 24-parameter motion model, five CompCor components derived from white-matter and cerebrospinal-fluid signals, and linear and quadratic temporal trends. This processing reduces contributions from residual head motion, physiological noise, and slow signal drift before connectivity is calculated. The resulting signals are then band-pass filtered between 0.01 and 0.10~Hz to retain the low-frequency fluctuations used for resting-state functional-connectivity analysis, while global signal regression is not applied under the selected \texttt{filt\_noglobal} strategy \citep{Craddocketal2013}. Following preprocessing, the functional data are reduced from voxel-level images to atlas-defined regional time series. For the primary analysis, the CC200 atlas divides the brain into 200 functionally defined regions, and the signal within each region is summarised across the scan to produce one time series per region and participant. These 200 regional time series form the common imaging input from which the correlation, static tangent-space, and windowed tangent-space connectivity representations are constructed. Although the common C-PAC workflow reduces variation caused by inconsistent preprocessing procedures, heterogeneity arising from acquisition conditions and participant composition remains part of the ABIDE-I multisite setting. These remaining sources of variation are addressed through quality-controlled cohort selection, site-aware validation, nuisance adjustment, and the sensitivity analyses described in the following sections.

\subsubsection{Primary and wider quality-control cohorts}
\label{subsubsec:cohorts}

The analysis cohorts are defined using the anatomical and functional image-quality ratings
provided with ABIDE-I. These ratings identify whether the available scans are suitable for
analysis and are accompanied by annotations for issues such as excessive head motion,
unsuccessful skull stripping, incomplete or distorted anatomical coverage, preprocessing
failure, image noise, missing imaging modalities, and regional signal dropout. For the
primary cohort, a participant is retained only when the required assessment from all three
raters is satisfied: rater~1 \reviewhl{is required to mark} the scan as acceptable; rater~2 \reviewhl{is required to rate} both the
anatomical and functional images as acceptable or borderline; and rater~3 \reviewhl{is required to rate} both
images as acceptable. Applying this criterion excludes 241 of the 1,112 ABIDE-I records.
Because the quality annotations \reviewhl{sometimes overlap}, an excluded scan \reviewhl{is potentially associated} with more
than one recorded issue. The resulting primary cohort contains 871 participants, including
403 autistic and 468 non-autistic participants. It comprises 727 male and 144 female
participants, with ages ranging from 6.5 to 58 years (mean 16.9 years; median 14.8 years).
Cohort membership is fixed before the primary analyses, and participant identifiers are
cross-checked against the phenotypic records to ensure that the same participants are used
throughout the prediction pipeline. The cohort-selection procedure and subsequent
time-series extraction are summarised in Figure~\ref{fig:overview}a.

Participant characteristics by diagnostic group are reported in
Table~\ref{tab:participants}. Welch's $t$-test is used to compare age and the three
intelligence quotient measures---full-scale, verbal, and performance IQ---between the
autistic and non-autistic groups, allowing for unequal group variances. Mean framewise
displacement (FD), defined in Section~\ref{subsubsec:phenotypic-variables}, is examined using both Welch's
$t$-test and the Mann--Whitney $U$ test. The former compares group means, whereas the latter
provides a rank-based comparison that is less sensitive to the right-skewed FD
distribution. Scan length is compared using the Mann--Whitney $U$ test, and sex
distributions are compared using the $\chi^2$ test with Yates' continuity correction.
Because cognitive ability and head motion \reviewhl{potentially influence} subsequent results, their effects
are examined separately in the cognitive-subgroup analysis of
Section~\ref{par:cognitive-sex-subgroups} and the motion-matched group-contrast analysis of
Section~\ref{subsubsec:motion-matching}.

A wider quality-control cohort is additionally defined to test whether the findings depend
on the more restrictive primary inclusion rule and to support comparison with previous ABIDE-I
studies. For this cohort, eligibility requires an acceptable functional-image rating from
rater~3, without applying the full three-rater criterion used for the primary cohort. This
definition identifies 1,035 eligible participants in the phenotypic records, of whom 1,032
have regional time series available for analysis. The 1,035-participant eligibility count
matches the cohort definition reported by \citet{Heinsfeldetal2018},
\citet{Eslamietal2019}, and \citet{AlmuqhimSaeed2021}. The wider cohort is therefore used
as a cohort-definition sensitivity analysis, allowing the primary results to be assessed
under a less restrictive quality-control rule while also providing a direct basis for
comparison with these previous reports.

\begin{table}[t]
\centering
\caption{Participant characteristics in the primary cohort (871 participants; 20 site labels from 17 acquisition sites). Values are mean (SD) unless stated; the number of participants with an available value is given in brackets where it is below the group total.}
\label{tab:participants}
\small
\begin{tabular}{@{}>{\raggedright\arraybackslash}p{0.36\linewidth}ccc@{}}
\toprule
Characteristic & Autistic (n = 403) & Non-autistic (n = 468) & p \\
\midrule
Age, years & 17.07 (7.96) & 16.84 (7.24) & 0.66 \\
Female, n (\%) & 54 (13.4) & 90 (19.2) & 0.027 \\
Full-scale IQ & 105.79 (17.06) [376] & 111.09 (12.08) [435] & $6.2\times10^{-7}$ \\
Verbal IQ & 104.89 (18.10) [342] & 111.57 (13.07) [379] & $2.8\times10^{-8}$ \\
Performance IQ & 105.69 (17.26) [344] & 107.91 (13.05) [391] & 0.052 \\
Mean framewise displacement (Section~\ref{subsubsec:phenotypic-variables}), mm & 0.126 (0.143) & 0.094 (0.082) & $8.2\times10^{-5}$ \\
Mean framewise displacement (Section~\ref{subsubsec:phenotypic-variables}), mm, median [IQR] & 0.088 [0.057, 0.144] & 0.068 [0.045, 0.109] & $8.9\times10^{-7}$ \\
Volumes per scan, median [range] & 176 [78, 316] & 176 [78, 316] & 0.77 \\
\bottomrule
\end{tabular}

\vspace{4pt}
\begin{minipage}{\linewidth}\footnotesize
p-values are from Welch's t-test for rows summarised by mean (SD), the Mann--Whitney U
test for rows summarised by median, and the $\chi^2$ test with Yates' continuity
correction for sex.
\end{minipage}
\end{table}

\subsubsection{Phenotypic variables}
\label{subsubsec:phenotypic-variables}

Phenotypic information is linked to each participant through the ABIDE-I identifiers and
includes age, sex, full-scale IQ, verbal IQ, performance IQ, handedness, and acquisition
site. These variables are used to characterise the cohort and to support the secondary,
subgroup, and nuisance-adjustment analyses described later. The primary ensemble itself is
restricted to imaging-derived connectivity features, so demographic and cognitive variables
do not contribute directly to its predictions. A separate phenotype model, described in
Section~\ref{par:phenotype-model}, instead uses age, sex, the three IQ measures, and handedness to
evaluate discrimination from non-imaging information. Acquisition site is retained for
site-based validation, nuisance adjustment, and sensitivity analyses. Head motion is represented separately by mean framewise displacement (FD), obtained from
the \texttt{func\_mean\_fd} \reviewhl{mean-FD quality metric} provided with the preprocessed release. This
measure follows Jenkinson's formulation and is calculated from the rigid-body realignment
parameters between successive functional volumes, with rotational components converted to
millimetres using an 80-mm head radius. The resulting mean FD
provides a single scan-level summary of head motion for each participant and is used in the
participant-characteristic comparisons and the subsequent motion-adjusted and
motion-matched analyses.

Missing phenotypic entries are identified before analysis, including values recorded as
\texttt{-9999} or left blank. When phenotypic variables are used in predictive or
harmonisation models, missing continuous values are replaced by the median calculated from
the corresponding outer training partition. This prevents information from held-out
participants from entering preprocessing. The same rule is used for the secondary
phenotype model in Section~\ref{par:phenotype-model} and for age in the ComBat analysis of
Section~\ref{par:combat-harmonisation}. In the ComBat model, sex is coded as 0 for male and 1 for female.
No age or sex values are missing in the primary cohort. For analyses fitted directly to a complete analysis cohort, required missing covariates are
median-imputed from that cohort. This applies to the nuisance-adjustment procedures in
Sections~\ref{par:selection-scope} and~\ref{subsubsec:connection-adjustment}, as well as to the preparation for motion
matching in Section~\ref{subsubsec:motion-matching}. In contrast, no imputation is applied in the
symptom-association, cognitive-subgroup, instrument-concordance, component--symptom, or
participant-characteristic analyses. Participants with missing values are excluded only
from the analysis that requires the corresponding variable.

\subsubsection{Atlas and signal quality}
\label{subsubsec:atlas-signal-quality}

The primary analyses use the Craddock 200-region functional parcellation (CC200)~\citep{Craddocketal2012}. All 200 regions are retained for prediction, yielding 19,900 unique pairwise connectivity features. A separate signal-quality check is applied before the group-difference analyses because Pearson correlation cannot be computed for a region with no temporal variation. First, any region with a temporal standard deviation of zero in more than 1\% of the 1,032 available CC200 time-series files is removed; four regions meet this criterion. Participants in the primary cohort are then excluded from the group-contrast analysis if any of the remaining regions contains a constant signal. This produces a group-contrast cohort of 863 participants and 196 regions, corresponding to 19,110 unique connections. To support system-level interpretation, each CC200 region is also assigned an anatomical label and a functional-system label. Yeo's seven cortical systems are mapped onto the CC200 space using nearest-neighbour interpolation~\citep{Yeoetal2011}, preserving a single system label for each voxel. Each region is assigned to the system with the greatest spatial overlap when that system covers more than 20\% of the region; regions that do not meet this threshold are assigned to a subcortical/cerebellar category. This procedure yields eight functional-system classes and 36 unique within- and between-system combinations, which are used to summarise connection-level group effects at the system level.

\subsubsection{Scan length}
\label{subsubsec:scan-length}

Scan length is considered explicitly in the dynamic-connectivity analyses because the number of available volumes differs substantially across participants and sites. In the primary cohort, scans contain between 78 and 316 volumes, with scan duration varying systematically by acquisition site. To keep the dynamic representations comparable, each participant contributes the same number of temporal windows regardless of total scan length. This prevents participants with longer scans from contributing more window-level estimates than those with shorter scans. Window length and temporal change are expressed in volumes rather than seconds because a harmonised repetition-time field is not available across the analysis files.

\subsection{Connectivity representations}
\label{subsec:connectivity-representations}

This subsection describes how each participant's regional rs-fMRI time series are
transformed into the three connectivity representations used by the primary ensemble:
shrinkage correlation, static tangent-space connectivity, and windowed tangent-space
connectivity. The first two describe connectivity over the full scan, whereas the third
summarises connectivity estimated from shorter temporal windows. Each representation is
converted into a 19,900-dimensional feature vector for the 200-region CC200 atlas and is
processed independently before classification. The corresponding feature construction,
covariance estimation, tangent-space projection, additional shrinkage, and windowing
procedures are detailed in Sections~\ref{subsubsec:connectivity-features}--\ref{subsubsec:windowed-representation}.

Figure~\ref{fig:representation} summarises this transformation. Whole-scan and window-level
covariance matrices are first estimated using Ledoit--Wolf shrinkage. Static tangent
connectivity is obtained by projecting the whole-scan covariance relative to a Riemannian
reference estimated from the corresponding training partition, while shrinkage correlation
is derived directly from the covariance estimate. For the windowed representation, each
scan is divided into six windows of 50 volumes, and the window covariances receive an
additional isotropic shrinkage of $r=0.15$ before projection to a training-derived dynamic
reference. The resulting window tangent matrices are averaged to form the mean windowed
tangent representation. All training-dependent reference matrices are re-estimated within
each outer fold and then applied unchanged to held-out participants.

\begin{figure}[pos=tbp]
\centering
\includegraphics[width=\linewidth]{fig_representation}
\caption{Construction of the connectivity representations for one example participant. (a) 1~Ledoit--Wolf covariance, 2~training-only Riemannian reference, 3~whitening and matrix logarithm, and 4~vectorisation; a shrinkage-correlation branch is also shown. (b) 5~Six windows of 50 volumes, each processed by steps 1--3 at a dynamic reference and then averaged. Blue lanes mark training-only quantities.}
\label{fig:representation}
\end{figure}
\subsubsection{Functional-connectivity features}
\label{subsubsec:connectivity-features}

Functional-connectivity features are derived from the regional rs-fMRI time series of each
participant. Each CC200 region is represented by a blood-oxygen-level-dependent signal
recorded across the scan, and connectivity is estimated from the statistical relationship
between pairs of these regional signals. With 200 regions, this produces 19,900 unique
pairwise connectivity features for each participant. Three representations of these
connections are used in the primary ensemble: shrinkage correlation, static tangent-space
connectivity, and mean windowed tangent-space connectivity. The static tangent
representation summarises connectivity across the full scan, whereas the windowed
representation captures the average connectivity pattern across shorter temporal segments.
These three representations are illustrated for one participant in
Figure~\ref{fig:overview}b. Each representation is constructed and standardised separately
so that its fitted reference and scaling parameters remain specific to that representation.

\subsubsection{Shrinkage-correlation features}
\label{subsubsec:shrinkage-correlation}

Shrinkage-correlation features are derived from each participant's Ledoit--Wolf covariance
estimate, defined in Section~\ref{subsubsec:covariance-estimation}. This estimate provides a regularised
description of how the regional signals vary together and is converted to a correlation
matrix by normalising each covariance by the standard deviations of the corresponding
region pair. The resulting coefficients therefore describe pairwise coupling on a common
scale, independent of regional signal variance. The strict upper triangle of the
correlation matrix is then vectorised into a single feature vector, while the diagonal is
omitted because all self-correlations equal one. This produces 19,900 correlation features
for the CC200 atlas. The same outer-fold partitions are used for this representation and
the two tangent-space representations, ensuring that all three ensemble members are
trained and evaluated on identical participant splits.

\subsubsection{Whole-scan covariance estimation}
\label{subsubsec:covariance-estimation}

Whole-scan covariance is estimated for each participant from all volumes of the regional
time series. For this purpose, $X_i \in \mathbb{R}^{T_i \times P}$ denotes the regional
time series for participant $i$, where $T_i$ is the number of volumes and $P$ is the
number of regions. The covariance is then estimated with Ledoit--Wolf shrinkage
\citep{LedoitWolf2004}, as given in Equation~\eqref{eq:ledoit-wolf-covariance}. In this estimator, the
covariance is a weighted average of the sample covariance and a scaled identity matrix:
\begin{equation}
\widehat{\Sigma}_i
=
(1-\rho_i)\widehat{S}_i
+
\rho_i \widehat{\mu}_i I_P,
\qquad
\widehat{\mu}_i
=
\frac{\operatorname{tr}(\widehat{S}_i)}{P}.
\label{eq:ledoit-wolf-covariance}
\end{equation}

In Equation~\eqref{eq:ledoit-wolf-covariance}, $\widehat{S}_i$ is the sample covariance, $I_P$ is the identity
matrix, and $\widehat{\mu}_i$ is the average of the regional variances. The shrinkage
weight $\rho_i$ is, in turn, selected by the Ledoit--Wolf estimator. Shrinkage is used
because the number of regions is large relative to the number of volumes, which makes the
sample covariance alone an unstable estimate.

\subsubsection{Tangent-space projection}
\label{subsubsec:tangent-projection}

Tangent-space projection converts each covariance matrix into a representation \reviewhl{suitable for}
used directly by standard linear classifiers. Because covariance matrices have a
non-Euclidean structure, they are projected relative to a reference matrix estimated from
the training data. This produces a locally linear feature space while preserving the
geometry of the original covariance matrices. The projection process is illustrated in
Figure~\ref{fig:representation}a.

The reference matrix is estimated separately for each outer fold as the affine-invariant
Riemannian mean of the training covariances \citep{Varoquauxetal2010,Dadietal2019}. As
given in Equation~\eqref{eq:riemannian-mean}, this mean is the positive-definite matrix $R$ with the
smallest summed squared distance to the training covariances:
\begin{equation}
\overline{\Sigma}_{\mathrm{T}}
=
\arg\min_{R \succ 0}
\sum_{i \in \mathrm{T}}
\left\|
\log\!\left(R^{-1/2}\widehat{\Sigma}_iR^{-1/2}\right)
\right\|_F^2.
\label{eq:riemannian-mean}
\end{equation}

In Equation~\eqref{eq:riemannian-mean}, $\mathrm{T}$ denotes the training set, and the Frobenius norm is used to quantify the affine-invariant distance between the reference matrix $R$ and each training covariance. The Riemannian mean is therefore adopted as the reference because it respects the geometry of the covariance space. It is estimated by fixed-point iteration with at most ten updates. Matrix logarithms and inverse square roots are computed by symmetric eigendecomposition, with eigenvalues clamped at $10^{-6}$ to maintain numerical stability.

Training and held-out covariances are then projected into the tangent space with the
training reference. As shown in Equation~\eqref{eq:tangent-projection}, each covariance is first multiplied
on both sides by the inverse square root of the reference. The matrix logarithm of this
product then gives the projected matrix $S_i$:
\begin{equation}
S_i
=
\log\!\left(
\overline{\Sigma}_{\mathrm{T}}^{-1/2}
\widehat{\Sigma}_i
\overline{\Sigma}_{\mathrm{T}}^{-1/2}
\right).
\label{eq:tangent-projection}
\end{equation}

The projected matrix in Equation~\eqref{eq:tangent-projection} is then vectorised by retaining only its
off-diagonal upper triangle, as given in Equation~\eqref{eq:vectorisation}. Only the upper triangle is
retained because the matrix is symmetric, so the entries below the diagonal repeat those
above it. Here, $\operatorname{vecu}$ stacks the entries above the diagonal into a vector:
\begin{equation}
x_i
=
\operatorname{vecu}(S_i)
\in
\mathbb{R}^{P(P-1)/2}.
\label{eq:vectorisation}
\end{equation}

The vectorisation retains only the strict upper triangle, so the $P$ diagonal entries are discarded and the off-diagonal entries are taken without the $\sqrt{2}$ weighting that a Frobenius-isometric map would apply. For $P=200$ regions this gives $P(P-1)/2=19{,}900$ features, which is the count used throughout and shown in Figure~\ref{fig:representation}. The mapping therefore preserves the ordering and relative magnitude of the off-diagonal tangent coefficients, but it is not an isometry of the full symmetric matrix, and the regional variance terms on the diagonal are omitted from the feature set.

For CC200, the vector in Equation~\eqref{eq:vectorisation} therefore contains 19,900 features. In
addition, the reference matrix is estimated from the training partition and applied
unchanged to held-out participants. As a result, no information from held-out
participants enters the tangent projection.

\subsubsection{Additional isotropic shrinkage}
\label{subsubsec:isotropic-shrinkage}

Additional isotropic shrinkage is used in the planned comparisons to further regularise the
Ledoit--Wolf covariance estimates. The primary static tangent representation uses the
Ledoit--Wolf estimate directly, corresponding to $r=0$. For comparison with the windowed
representation, an additional shrinkage level of $r=0.15$ is applied using
Equation~\eqref{eq:isotropic-shrinkage}:
\begin{equation}
\widehat{\Sigma}_i(r)
=
(1-r)\widehat{\Sigma}_i
+
r\frac{\operatorname{tr}(\widehat{\Sigma}_i)}{P}I_P,
\qquad 0 \leq r < 1.
\label{eq:isotropic-shrinkage}
\end{equation}
A separate nested sensitivity analysis selects $r$ from
$\{0,0.10,0.20,0.30,0.40,0.50\}$ within each outer training partition. Classifier
regularisation is reselected for every candidate value of $r$, so each shrinkage setting
is evaluated with its corresponding tuned classifier while the held-out data remain
excluded from parameter selection.

\subsubsection{Windowed tangent-space representation}
\label{subsubsec:windowed-representation}

The windowed tangent-space representation is built from covariance matrices estimated
within the windows of each scan. In this approach, window-level covariance matrices are used
to represent dynamic connectivity
\citep{Hutchisonetal2013,Allenetal2014,LeonardiVanDeVille2015,Pretietal2017}. From these matrices, the mean windowed tangent representation, hereafter termed the dynamic
mean, forms the third member of the primary ensemble. In contrast, dispersion, velocity, and alternative block combinations, that is,
combinations of the mean, dispersion, and velocity blocks, are secondary analyses.
Dispersion and velocity are, in turn, defined below.
The placement of the six windows and the formation of the dynamic mean from the window
tangent matrices are illustrated in Figure~\ref{fig:representation}b.

For every participant, $K=6$ windows of $L=50$ volumes are extracted. The values $K=6$
and $L=50$ are used because every scan contains at least 78 volumes, so that each window
spans 50 volumes. The window positions are given in Equation~\eqref{eq:window-positions}, which spaces the windows
evenly from the start to the end of each scan. Specifically, for a scan with $T_i$
volumes, $\Delta_i$ is the nominal stride between consecutive window starts, and $s_{ik}$
is the start of window $k$:
\begin{equation}
\Delta_i
=
\frac{T_i-L}{K-1},
\qquad
s_{ik}
=
\operatorname{round}(k\Delta_i),
\qquad
k=0,\ldots,K-1.
\label{eq:window-positions}
\end{equation}

The number of windows is fixed across participants to keep the dynamic summaries comparable
despite differences in scan length across sites. Shorter scans therefore contain more
overlap between windows, while longer scans \reviewhl{sometimes contain} small gaps, and each window start
is rounded to the nearest acquired volume. For every window, covariance is estimated using
Ledoit--Wolf shrinkage with an additional isotropic shrinkage of $r=0.15$ to improve
numerical stability. Adaptive diagonal jitter is applied when required during
eigendecomposition. The dynamic training reference is then constructed by first averaging
the six window covariances for each training participant and subsequently taking the
Riemannian mean across participants. This procedure gives each training participant equal
weight in the reference.

Three summaries are then computed from the projected windows of each participant. After
tangent projection, participant $i$ is represented by the window vectors
$\xi_i(1),\ldots,\xi_i(K)$, which form a trajectory across the scan. Each summary is
computed elementwise, that is, separately for every connection. Specifically, the mean
$m_i$, the dispersion $d_i$, and the velocity $v_i$ are given in Equations~\eqref{eq:window-mean}, %tareq change here
\eqref{eq:window-dispersion}, and \eqref{eq:window-velocity}, respectively:
\begin{equation}
m_i
=
\frac{1}{K}\sum_{k=1}^{K}\xi_i(k),
\label{eq:window-mean}
\end{equation}
\begin{equation}
d_i
=
\left[
\frac{1}{K}
\sum_{k=1}^{K}
\left(\xi_i(k)-m_i\right)^2
\right]^{1/2},
\label{eq:window-dispersion} %tareq change here
\end{equation}
\begin{equation}
v_i
=
\frac{1}{(K-1)\Delta_i}
\sum_{k=1}^{K-1}
\left|\xi_i(k+1)-\xi_i(k)\right|.
\label{eq:window-velocity} %tareq change here
\end{equation}

The three summaries capture complementary properties of the window trajectory. The mean in
Equation~\eqref{eq:window-mean} represents the average connectivity pattern across windows, while the
dispersion in Equation~\eqref{eq:window-dispersion} quantifies within-scan variability around this mean
using the denominator $K$. The velocity in Equation~\eqref{eq:window-velocity} measures the average
absolute change between consecutive windows. It is normalised by the stride from
Equation~\eqref{eq:window-positions}, allowing temporal change to be compared across scans with different
window spacing.

\subsection{Prediction model}
\label{subsec:prediction-model}

This subsection describes how the three connectivity representations are converted into
participant-level predictions. Each representation is standardised and classified
separately using regularised logistic regression, after which the three probabilities are
combined to form the primary ensemble prediction. Sections~\ref{subsubsec:feature-scaling}--\ref{subsubsec:probability-ensemble}
describe scaling, classification, and probability averaging, while
Section~\ref{subsubsec:rank-fusion} presents the secondary rank-fusion analysis.

\subsubsection{Feature scaling}
\label{subsubsec:feature-scaling}

Before classification, each feature block, that is, the features of one representation, is
standardised. To this end, every feature is
centred on its mean and divided by its standard deviation. Both statistics are estimated
from the outer training partition, as shown in Equation~\eqref{eq:feature-scaling}:
\begin{equation}
z_{ij}
=
\frac{x_{ij}-\mu_{j,\mathrm{T}}}
{\sigma_{j,\mathrm{T}}}.
\label{eq:feature-scaling}
\end{equation}
The per-representation scaling, the regularised logistic models, and the probability
ensemble are summarised in Figure~\ref{fig:overview}c.

In Equation~\eqref{eq:feature-scaling}, $x_{ij}$ is feature $j$ of participant $i$, and
$\mu_{j,\mathrm{T}}$ and $\sigma_{j,\mathrm{T}}$ are the mean and standard deviation of
that feature over the training indices $\mathrm{T}$. The transformation is fitted on the
training partition and applied unchanged to the held-out partition. Because both
statistics come from the training partition alone, no information from held-out
participants enters the scaling. Standardisation, in turn, places all features on a common
scale before the penalised classifier is fitted. In addition, separate scalers are fitted
for the correlation, static tangent, and dynamic tangent members, so that each member is
standardised with its own training statistics.

\subsubsection{Regularised logistic regression}
\label{subsubsec:logistic-regression}

Logistic regression with an $\ell_2$ penalty is used for each primary member. Logistic
regression is a linear model that converts a weighted sum of the features into a
probability. The $\ell_2$ penalty, in turn, shrinks the weights towards zero. The penalty
is used because the number of connectivity features exceeds the number of training
participants. For participant $i$, the predicted probability is given by
Equation~\eqref{eq:logistic-probability}:
\begin{equation}
p_i
=
\sigma(w^\top z_i+b),
\qquad
\sigma(a)
=
\frac{1}{1+\exp(-a)}.
\label{eq:logistic-probability}
\end{equation}

In Equation~\eqref{eq:logistic-probability}, $w$ and $b$ denote the model weights and intercept, $z_i$ is
the standardised feature vector, and $\sigma$ is the logistic function. Logistic regression
is fitted in scikit-learn \citep{Pedregosaetal2011} using the \texttt{liblinear} solver
with a maximum of 2,000 iterations. The regularisation parameter $C$ is selected separately
for each representation from $\{0.001,0.01,0.1,1.0\}$ using three-fold stratified
cross-validation within the outer training partition, with AUROC used for selection. Class
weights are kept uniform. For the primary ensemble, the tangent reference and feature
scaler are fitted once on the complete outer training partition and reused during the inner
search for $C$. The outer held-out partition remains excluded from both transformation
fitting and model selection.

\subsubsection{Probability ensemble}
\label{subsubsec:probability-ensemble}

In the primary model, the probabilities from whole-scan correlation, whole-scan tangent
connectivity, and mean windowed tangent connectivity are averaged. For member
probabilities $p_i^{(1)},p_i^{(2)},p_i^{(3)}$, the ensemble probability is given by
Equation~\eqref{eq:probability-ensemble}:
\begin{equation}
p_i^{\mathrm{ens}}
=
\frac{1}{3}
\sum_{m=1}^{3}p_i^{(m)}.
\label{eq:probability-ensemble}
\end{equation}

The members in Equation~\eqref{eq:probability-ensemble} are given equal weight. Moreover, ensemble
membership and weighting are fixed across outer folds. Equal weighting is used because it
avoids estimation of additional fusion parameters from the training sample. Finally, the
ensemble probability is converted into a class decision with a default class threshold of
0.5.

\subsubsection{Secondary rank fusion}
\label{subsubsec:rank-fusion}

As a secondary analysis, the three model outputs are combined using their relative ranks rather than their predicted probabilities. For each held-out participant, a member score is located within the distribution of scores produced for that member's training participants. This converts the score to a value between 0 and 1 that reflects its relative position within the training distribution. When scores are tied, the held-out value is placed after the equal training scores. The three rank-based scores are then averaged to obtain the fused score. This approach places all ensemble members on a common scale while avoiding dependence on their original probability ranges. The fused value is therefore interpreted as an ordering score, not as a probability, and each held-out participant is ranked only against training scores.


\subsection{Validation protocol}
\label{subsec:validation-protocol}

Two outer validation schemes are used to assess model performance: pooled ten-fold cross-validation and leave-one-site-out validation. Within each outer training partition, classifier regularisation is selected through inner cross-validation, resulting in a nested evaluation design. All training-dependent transformations and model parameters are estimated from the outer training data only, keeping the held-out partition separate from model fitting and model selection. Outer validation and fold-level fitting is discussed in Section \ref{subsubsec:outer-validation}. Performance measures and fold-level aggregation are presented in Section~\ref{subsubsec:performance-measures}, while interval estimation is described in Section~\ref{subsubsec:interval-construction}.

\subsubsection{Outer validation and fold-level fitting}
\label{subsubsec:outer-validation}
Model evaluation is performed using both pooled ten-fold cross-validation and leave-one-site-out validation. In pooled ten-fold validation, the primary cohort is divided into ten stratified folds with similar diagnostic proportions, using shuffling with random seed 42. Each fold is held out once, so every participant contributes one test prediction and performance reflects generalisation within the combined ABIDE-I site mixture. In leave-one-site-out validation, each of the 20 recorded site labels is held out in turn, while participants from the remaining sites are used for training. This provides a stricter assessment of transfer to an ABIDE-I site that is absent from model fitting. Within every outer training partition, stratified inner folds are used to select the regularisation parameter $C$. Model fitting then follows the same sequence for both validation schemes: training and held-out indices are first defined, the tangent reference is estimated from the training covariances, and both partitions are projected using this fixed reference. Feature scaling is subsequently fitted on the training features and applied unchanged to the held-out data before classification. This procedure keeps all fitted transformations and model-selection steps within the training partition. The two outer validation schemes are summarised in Figure~\ref{fig:overview}c, while the nested fitting procedure is illustrated in Figure~\ref{fig:designs}a.

The scope of this nesting is stated precisely, because it is complete at the outer level only.
The tangent reference and the feature scaler are estimated once on each outer-training
partition and are then held fixed while the regularisation parameter is selected by inner
three-fold cross-validation. The inner-validation participants therefore contribute to the
reference and to the scaling statistics that are applied to them, although their diagnostic
labels are never used at any point. The consequence is confined to the selection of a single
hyperparameter within the training partition, and every outer held-out participant remains
excluded from all fitted quantities. The design is accordingly described as nested at the
outer level rather than fully nested, and a fully nested variant in which both transformations
are refitted within every inner split is identified as outstanding work.

Figure~\ref{fig:designs} summarises the main validation and comparative designs used in the study. Panel~(a) shows the nested fitting procedure: Ledoit--Wolf covariances are computed before the outer split, while tangent references, scaling parameters, and classifier tuning are restricted to the outer training partition. The selected classifier is then refitted, and the three member probabilities are averaged to form the ensemble prediction. Panel~(b) compares connection- and region-level feature selection, including both training-fold selection and deliberately global label-informed selection. Panel~(c) presents the group-contrast analysis, where connectivity effects are estimated in the full cohort and in within-site samples matched on framewise displacement, followed by within-site permutation testing and BH-FDR correction. Panel~(d) summarises the cognitive analyses, including comparisons across autistic IQ bands, matched comparisons with the non-autistic group, and within-group score--IQ associations. Panel~(e) compares linear and tree-based classifiers under the same static tangent representation, scaling procedure, and outer folds, with only the $\ell_2$-regularised logistic model receiving hyperparameter tuning. The arrows and colour conventions follow Figure~\ref{fig:overview}, while the remaining markers and matrices are schematic.


\begin{figure}[pos=tbp]
\centering
\includegraphics[width=\linewidth,height=0.78\textheight,keepaspectratio]{fig_designs}
\caption{Evaluation and analysis designs. (a) Nested cross-validation in the outer training partition: \ding{202}~reference, \ding{203}~projection, \ding{204}~scaling, \ding{205}~inner-CV tuning, \ding{206}~refitting. (b) Feature selection within training folds versus with all labels. (c) Motion-matched group contrast. (d) IQ designs: (i)~quartiles, (ii)~fixed comparison, (iii)~matching, (iv)~score--IQ correlation. (e) Classifier families.}
\label{fig:designs}
\end{figure}

The classifier is then tuned and refitted on the transformed training data. To this end,
$C$ is selected by inner cross-validation on the transformed training matrix. The selected
classifier is then refitted to the complete outer training partition, so that the final
model uses every training participant. Next, probabilities are predicted for the held-out
partition. Finally, the three member probabilities are combined.

Per-participant covariance estimation is handled separately from these fold-specific
steps. Each covariance estimate depends only on that participant's time series. Therefore,
it is cached before fold-specific fitting without using information from other
participants. Shared quantities, in contrast, are estimated from several participants. For
this reason, shared references, scalers, feature selectors, and classifier parameters are
estimated at the training level specified for each analysis.

\subsubsection{Performance measures and aggregation}
\label{subsubsec:performance-measures}

Discrimination is evaluated primarily using AUROC and the area under the precision--recall
curve, both calculated from continuous prediction scores. Threshold-based performance is
also reported at 0.5 using accuracy, balanced accuracy, sensitivity, specificity, F1 score,
and Matthews correlation coefficient, with autistic participants treated as the positive
class. Balanced accuracy and Matthews correlation coefficient are included to provide
measures that are less affected by unequal group sizes.

Performance is summarised in two ways across the outer folds. Participant-pooled metrics
are calculated from all concatenated out-of-fold predictions, while unweighted fold means
and standard deviations give equal importance to each fold. Both summaries are retained
because site sizes vary substantially in leave-one-site-out validation. Planned comparison
tables report fold-level AUROC and accuracy, whereas pooled out-of-fold predictions are
used for calibration, paired prediction tests, and the main interval estimates.

\subsubsection{Interval construction}
\label{subsubsec:interval-construction}

Intervals are reported for the main performance measures. AUROC intervals are obtained from
the DeLong variance, while Wilson intervals are used for accuracy,
sensitivity, and specificity. These intervals are conditional on the observed out-of-fold
predictions and therefore do not include uncertainty from repeating the complete modelling
procedure.

To summarise variation across acquisition sites, site-specific AUROCs from leave-one-site-out
validation are combined using a DerSimonian--Laird random-effects model. Only site labels containing both diagnostic groups are included,
leaving $J$ sites for analysis. Each site is weighted by
$1/(V_j+\hat{\tau}^2)$, where $V_j$ is its DeLong variance and $\hat{\tau}^2$ is the
estimated between-site variance. Between-site heterogeneity is summarised using Cochran's
$Q$ and $I^2$, and the expected range of AUROC at a new site is estimated with the 95\%
prediction interval in Equation~\eqref{eq:prediction-interval}:
\begin{equation}
\hat{\mu}
\pm
t_{0.975,\,J-2}
\sqrt{\hat{\tau}^2+\operatorname{SE}(\hat{\mu})^2}.
\label{eq:prediction-interval}
\end{equation}

Here, $\hat{\mu}$ is the random-effects mean AUROC and
$\operatorname{SE}(\hat{\mu})$ is its standard error. By incorporating both uncertainty in
the pooled mean and estimated between-site variation, the interval \reviewhl{summarises descriptive variation among the ABIDE-I site estimates}.

The site-specific estimates are not independent studies, because each is produced by a model
trained on the remaining 19 sites, so any two training sets overlap in almost all of their
participants. The random-effects model does not represent this dependence, and the
between-site variance is therefore understated. For this reason the random-effects mean and
the interval are reported as descriptive summaries of variation among the ABIDE-I site
estimates. Neither quantity is interpreted as the performance expected at a new acquisition
site, and the site-clustered bootstrap of Section~\ref{subsubsec:paired-comparisons} is used
where an interval on overall performance is required.



\subsection{Probability assessment and selective prediction}
\label{subsec:probability-assessment}
%[need one intro sentence here]
This subsection examines the reliability of the ensemble probabilities through calibration, risk--coverage analysis, and confidence-based selective prediction with finite-sample error control.
\subsubsection{Secondary probability assessment}
\label{subsubsec:calibration-assessment}

Beyond discrimination, the ensemble probabilities are evaluated for calibration and later used in selective prediction. The calibration analysis examines the native probabilities from Equation~\eqref{eq:probability-ensemble} without any additional recalibration. Their reliability is assessed by comparing predicted probabilities with the observed proportion of autistic participants. The probabilities are divided into ten equal-width bins, and the mean prediction in each bin is compared with the corresponding observed outcome frequency. These bin-level differences are then combined into the expected calibration error (ECE), weighted by the number of participants in each bin \citep{Naeinietal2015,Guoetal2017}, as defined in Equation~\eqref{eq:calibration-error}:

Because a single binned summary does not establish that probabilities are calibrated, three
further measures are reported. The Brier score is computed as the mean squared difference
between the predicted probability and the outcome, so that discrimination and calibration are
summarised together. The calibration slope is obtained by regressing the outcome on the
logit of the predicted probability, where a slope near one indicates that the probabilities
are neither over- nor under-dispersed. Calibration-in-the-large is obtained from a model
containing only an intercept, with the logit of the predicted probability entered as a fixed
offset, so that a value near zero indicates that the average predicted risk matches the
observed rate. Each of these measures is reported with the site-clustered bootstrap interval
of Section~\ref{subsubsec:paired-comparisons}. Sensitivity to binning is examined by
recomputing the expected calibration error with 5, 10 and 20 equal-width bins, and the
smallest bin count is reported for each choice.

\begin{equation}
\operatorname{ECE}
=
\sum_{b=1}^{10}
\frac{|\mathrm{B}_b|}{n}
\left|
\frac{1}{|\mathrm{B}_b|}
\sum_{i\in\mathrm{B}_b}p_i^{\mathrm{ens}}
-
\frac{1}{|\mathrm{B}_b|}
\sum_{i\in\mathrm{B}_b}y_i
\right|.
\label{eq:calibration-error}
\end{equation}

In Equation~\eqref{eq:calibration-error}, $\mathrm{B}_b$ denotes the set of participants in bin $b$,
$n$ denotes the total number of participants, and $y_i=1$ denotes an autistic participant.
In this sum, empty bins contribute no term. Because ECE is computed from the probability of
the autistic class, it assesses positive-class probability calibration. In addition, the
same bin-level quantities are displayed in reliability diagrams under both validation
protocols. In such a diagram, the mean predicted probability of each bin is plotted against
the observed proportion of autistic participants in that bin. These diagrams are included
because they show in which probability range any disagreement occurs, whereas a single ECE
value does not.

\subsubsection{Exploratory risk--coverage analysis}
\label{subsubsec:risk-coverage}

In the exploratory risk--coverage analysis, the change in error rate is examined as
predictions are retained only for the most confident participants. In this analysis,
coverage is the fraction of participants whose predictions are retained. The analysis is
performed because it shows whether higher model confidence corresponds to fewer incorrect
decisions. To this end, prediction confidence is defined in Equation~\eqref{eq:prediction-confidence} as the
larger of the ensemble probability and its complement:
\begin{equation}
c_i
=
\max\left(p_i^{\mathrm{ens}},1-p_i^{\mathrm{ens}}\right).
\label{eq:prediction-confidence}
\end{equation}

Participants are then ordered by the confidence in Equation~\eqref{eq:prediction-confidence}. Next, the number
of selected participants, accuracy, AUROC, and selective risk are recorded at coverage
levels from 10\% to 100\%. Here, selective risk is the fraction of incorrect decisions
among the selected participants \citep{GeifmanElYaniv2017}. Specifically, at each requested
coverage, the highest-confidence observations in the complete out-of-fold prediction set
are retained. However, AUROC is undefined when a retained subset contains one diagnostic
class, because AUROC requires participants from both classes.

\subsubsection{Exploratory finite-sample threshold analysis}
\label{subsubsec:threshold-analysis}

In the exploratory threshold analysis, a confidence threshold $\theta$ is selected from
labelled calibration data. In this context, calibration data are participants with known
diagnostic labels who are used to choose the threshold. For this selection, the candidate
thresholds are the distinct observed confidence values, and a participant is selected when
$c_i\geq\theta$. For each candidate $\theta$, $k(\theta)$ denotes the number of selected
participants, and $e(\theta)$ denotes the number of errors among them. The error rate of
the selected participants is then bounded with the one-sided Clopper--Pearson upper bound
\citep{ClopperPearson1934}, as given in Equation~\eqref{eq:clopper-pearson-bound}:
\begin{equation}
U(\theta;\delta)
=
\operatorname{BetaInv}
\left(
1-\delta;
e(\theta)+1,
k(\theta)-e(\theta)
\right).
\label{eq:clopper-pearson-bound}
\end{equation}

In Equation~\eqref{eq:clopper-pearson-bound}, $\operatorname{BetaInv}$ is the beta-distribution quantile
function and $1-\delta$ sets the confidence level. The upper confidence bound is used
instead of the observed error rate so that threshold selection accounts for sampling
uncertainty in the calibration data. Candidate thresholds are evaluated in a fixed order,
from the most restrictive to the most permissive. The search stops when the upper bound
first exceeds the target risk $\alpha$, and the last passing threshold is retained.
Candidates selecting fewer than ten participants are skipped
\citep{Angelopoulosetal2021,AngelopoulosBates2023}. The fixed ordering allows each
candidate to be tested at level $\delta$ without an additional correction across
thresholds.

\reviewhl{[AUTHORS TO VALIDATE: Confirm against the cited learn-then-test procedure that the candidate order, stopping rule, reused fallback threshold, and unstratified half-split evaluation satisfy the claimed finite-sample guarantee. State the precise guarantee and its exchangeability assumptions; otherwise describe this as an empirical threshold-selection experiment.]}

Two boundary cases are handled explicitly. If the first eligible candidate already
exceeds the target risk, the threshold is set to 1.01, which gives zero coverage because
confidence scores cannot exceed one. If every selected case is an error, the upper bound
is set to one because the beta quantile in Equation~\eqref{eq:clopper-pearson-bound} is otherwise undefined.
For the ensemble analysis, $\alpha=0.25$, $\delta=0.10$, and a minimum selected sample of
ten are used, corresponding to a 25\% target error with a 90\% upper confidence bound.
Two thresholding strategies are evaluated. Marginal thresholding uses one threshold for
the full calibration sample, whereas site-conditional thresholding estimates a separate
threshold for each site. When a site's calibration subset contains fewer than ten
participants, the marginal threshold is used instead.

Threshold performance is assessed over 200 random, unstratified half-splits. In each
repetition, one half of the data is used to select the threshold and the other half is used
to measure realised coverage and selective risk. This keeps threshold evaluation separate
from threshold selection. Repetitions with undefined selective risk are excluded from the
summary. The reported results include the mean, median, and 90th-percentile risk, mean
coverage, and the proportion of valid repetitions that satisfy the target risk. Repeated
splitting is used to capture variation across different calibration samples rather than
relying on a single partition.



\subsection{Planned comparisons}
\label{subsec:planned-comparisons}
This subsection presents the planned primary, secondary, and exploratory comparisons built on the same prediction and validation framework. These analyses examine how performance and interpretation vary with connectivity representation, classifier choice, atlas and feature-selection strategy, phenotypic and site effects, participant subgroups, data availability, and model interpretability. 

\subsubsection{Primary descriptive comparison}
\label{subsubsec:primary-comparisons}

\paragraph{\textbf{Static and windowed connectivity.}}
\label{par:static-windowed}
The planned comparisons use the three-member prediction model under both pooled ten-fold and leave-one-site-out validation, with each analysis classified by its status in Section~\ref{subsubsec:analysis-status}. The first primary descriptive comparison evaluates static tangent connectivity against the dynamic mean while controlling the level of additional shrinkage. Because the windowed representation uses $r=0.15$, static tangent connectivity is evaluated at both $r=0$ and $r=0.15$, and the dynamic mean is compared directly with the $r=0.15$ static representation. Other representation-specific steps, including window sampling, covariance aggregation, and reference construction, remain unchanged. Static tangent performance is also examined across the shrinkage grid ${0,0.05,0.10,0.15,0.20,0.30,0.40,0.50,0.65,0.80}$. The sweep is first evaluated up to 0.30 and is later extended to 0.80 to examine stronger shrinkage, without changing the fixed $r=0.15$ comparison. In a separate nested analysis, the shrinkage level is selected within each outer training partition. The fixed and data-selected analyses are therefore reported separately, since one uses a predefined shrinkage level while the other estimates it from the training data.

\paragraph{\textbf{Feature-selection scope.}}
\label{par:selection-scope}
This primary descriptive comparison examines how model performance changes when fitted steps are allowed to access data beyond the outer training fold. Seven data-access conditions are evaluated using the 19,900 static tangent features from the primary cohort, while the same participants and outer folds are retained throughout. The reference condition fits the tangent reference, scaler, and classifier within each outer training fold. The global-reference condition estimates the tangent reference from all participants but retains fold-specific scaling, while the global-scaling condition fits both the reference and scaler globally. In the global-selection condition, the tangent reference is fitted globally, scaling remains fold-specific, and the 2,000 strongest connections are selected with an F test using all diagnostic labels. These features represent approximately 10\% of all connections. A held-out-tuning condition selects $C$ using the held-out fold, whereas a site-feature condition adds one-hot site indicators to the otherwise fold-specific pipeline. Finally, the combined condition applies global tangent projection, global scaling, global label-informed feature selection, held-out $C$ selection, and site indicators together. Here, ``global'' denotes any step fitted using all participants, including those in the held-out fold. The main connection-level conditions are illustrated in Figure~\ref{fig:designs}b.

The same question is also examined at the region level. This analysis uses the group-contrast cohort of 863 participants and 196 regions, with additional shrinkage fixed at $r=0.30$. The nuisance model is first fitted to the complete analysis cohort, after which each region is ranked by the average absolute group difference across its incident connections. Region ranking is performed either with all diagnostic labels or with labels restricted to the outer training fold, allowing the effect of label access to be isolated. Models are then fitted using the top 10, 20, 30, 50, or 100 ranked regions, while a model containing all retained regions serves as the reference condition.

In the published training-fold condition the regions are ranked inside each fold, but the nuisance adjustment used for that ranking is fitted once on the complete analysis cohort, so the held-out participants contribute to the residuals from which the ranking is derived. The condition is therefore repeated with the nuisance coefficients estimated from the outer-training rows alone and applied to those rows only, so that every data-dependent step of the selection sees training participants and nothing else. The folds, the connectivity features, the tangent projection, the feature scaling and the classifier are unchanged, so any difference is attributable to the nuisance fit.

\paragraph{\textbf{Pooled and within-participant symptom associations.}}
\label{par:symptom-associations}
This primary descriptive comparison examines whether out-of-fold ensemble probabilities are associated with six symptom measures. These include the ADOS total and calibrated severity scores, the ADI-R social, verbal-communication, and restricted/repetitive-behaviour domain scores, and the SRS total \citep{Lordetal1994,Lordetal2000}. Each analysis uses only participants with an available score for the corresponding instrument, and missing symptom values are not imputed.

Associations are measured using Spearman correlation and are evaluated at two levels. The pooled analysis includes all participants with an available symptom score, whereas the within-participant analysis is restricted to autistic participants. This second analysis is included because an association observed across both diagnostic groups \reviewhl{partly reflects} the difference between the groups rather than variation within autistic participants. Both analyses are performed using out-of-fold predictions from pooled ten-fold and leave-one-site-out validation.



\subsubsection{Secondary comparison}
\label{subsubsec:secondary-comparisons}

\paragraph{\textbf{Dynamic trajectory blocks.}}
\label{par:trajectory-blocks}
In this secondary comparison, the dynamic mean is evaluated both alone and together with dispersion and velocity. Two combination strategies are examined. In feature concatenation, the mean is combined with dispersion, with velocity, or with both, and a single classifier is fitted to each joined feature vector. In late fusion, separate classifiers are fitted to the individual trajectory blocks and their predicted probabilities are averaged with equal weights. These approaches are compared because concatenation applies one regularisation penalty across all included features, whereas late fusion allows each block to be modelled separately. Standalone performance is reported for the dynamic mean, while dispersion and velocity are evaluated only as part of the combined conditions.

\paragraph{\textbf{Classifier family.}}
\label{par:classifier-family}
Eight classifier configurations are compared using the same static tangent features, feature scaling, and outer validation folds. This common setup ensures that performance differences are not caused by changes in the input features or participant partitions. Five linear or low-rank models are evaluated: nested $\ell_2$-regularised logistic regression; ridge classification with penalty strength 100; linear support-vector classification with $C=0.01$; shrinkage linear discriminant analysis after reduction to 200 principal components; and a rank-20 two-stage low-rank spectral classifier, which is specified below. Three tree ensembles are also included. Random forest uses 300 trees with unrestricted depth, XGBoost uses 300 trees with maximum depth 4, learning rate 0.1, row subsampling 0.8, and column subsampling 0.3, and CatBoost uses 300 trees with depth 4, learning rate 0.1, and feature subsampling 0.3. The linear and tree-based configurations are summarised in Figure~\ref{fig:designs}e.

Hyperparameter tuning is applied only to the logistic model because it is the classifier used in the primary ensemble. Its $C$ value is selected from $\{0.001,0.01,0.1,1.0\}$ using three-fold stratified cross-validation within each outer training partition. All other classifiers are evaluated with fixed configurations \citep{Prokhorenkovaetal2018}. The analysis is therefore treated as a secondary comparison of the specified configurations rather than a comparison of the best attainable performance of each classifier family. Before threshold-based metrics are calculated, all model outputs are placed on a common zero-to-one scale. For classifiers without \texttt{predict\_proba}, the decision score is transformed with a logistic function, allowing the same decision threshold to be applied across classifiers.

\paragraph{\textbf{Phenotype model.}}
\label{par:phenotype-model}
In this secondary comparison, a phenotype-only model is fitted using age, sex, full-scale IQ, verbal IQ, performance IQ, and handedness. Missing values are imputed using medians estimated from the outer training partition, after which the variables are standardised and entered into logistic regression. The resulting phenotype probability is evaluated both on its own and as an additional member of the imaging ensemble. This allows the contribution of phenotype information to be examined separately from the discrimination provided by the imaging models.

\paragraph{\textbf{Atlas and fusion.}}
\label{par:atlas-fusion}
This secondary comparison examines whether connectivity described by more than one brain parcellation improves discrimination. In addition to CC200, the automated anatomical labelling atlas \citep{TzourioMazoyeretal2002} and the Harvard--Oxford probabilistic atlas are used, both taken from the regional time series distributed with the ABIDE Preprocessed Initiative \citep{Craddocketal2013}, where the Harvard--Oxford derivative comprises 111 regions, because these two atlases divide the brain on anatomical rather than functional grounds and \reviewhl{therefore potentially carry} complementary information.  Static tangent features are computed separately for each atlas, and the probabilities of the corresponding models are averaged for every combination of two and three atlases. The same outer folds, feature scaling, and classifier are used throughout, so that any difference in performance \reviewhl{is attributable} to the atlas information alone.

\paragraph{\textbf{ComBat harmonisation.}}
\label{par:combat-harmonisation}

In this secondary comparison, ComBat harmonisation is used to reduce site-related shifts in feature means and variances \citep{Johnsonetal2007,Fortinetal2018}. Site is treated as the batch variable, while age and sex are retained as covariates and diagnosis is excluded so that harmonisation does not depend on the prediction target. ComBat is fitted separately within each outer fold after feature scaling. The training features and the complete unlabelled held-out batch are supplied jointly so that a batch effect \reviewhl{is estimated} for a site that is absent from the training data. \reviewhl{This is a transductive batch-level analysis because it requires simultaneous access to the full unlabelled held-out site; it is not an inductive test for a single participant arriving from an unseen site.} If age or sex were missing, age would be replaced by the training median and sex would be assigned to the male category, although these rules are not required in the primary cohort. When ComBat-transformed features are unavailable, the scaled features are retained under a prespecified fallback so that every outer fold still contributes predictions.


\paragraph{\textbf{Low-rank spectral approximation.}}
\label{par:low-rank-spectral}
In this secondary analysis, a linear classifier is approximated with a small number of
spectral components of its weight matrix, that is, with a low-rank spectral approximation.
To this end, each standardised static tangent feature vector is represented as the upper
triangle of a symmetric matrix $\widetilde{S}_i$. In the same way, a dense linear
coefficient vector, that is, one with a non-zero weight for every feature, is mapped to a
symmetric weight matrix $W$. The off-diagonal entries of $W$ are divided by two, so that
the matrix inner product matches the original vector dot product.

The weight matrix is then decomposed into its eigenvectors and eigenvalues, as given in
Equation~\eqref{eq:weight-eigendecomposition}:
\begin{equation}
W
=
U\Lambda U^\top.
\label{eq:weight-eigendecomposition}
\end{equation}

In Equation~\eqref{eq:weight-eigendecomposition}, the columns of $U$ are the eigenvectors of $W$, and the diagonal
matrix $\Lambda$ holds the eigenvalues. Each eigenvector, together with its eigenvalue,
forms one component. The components are ordered by the absolute value of their eigenvalues.
This ordering places first the components that contribute most to the weight matrix. For
rank $q$, that is, when $q$ components are retained, the first $q$ components then define
the participant scores in Equation~\eqref{eq:component-scores}:
\begin{equation}
z_{ik}
=
u_k^\top \widetilde{S}_i u_k,
\qquad k=1,\ldots,q.
\label{eq:component-scores}
\end{equation}

In Equation~\eqref{eq:component-scores}, $u_k$ is the $k$th column of $U$, and $z_{ik}$ is the score of
participant $i$ on component $k$. The scores in Equation~\eqref{eq:component-scores} are obtained from a
dense stage-one logistic model with $C=0.01$. Next, a second logistic regression with
$C=1.0$ is fitted to the component scores. This two-stage design is used so that the final
classifier operates on the $q$ component scores alone. The ranks $\{1,2,5,10,20,30,50,100\}$
are evaluated on the same folds, so that differences between ranks do not reflect different
data partitions. In a separate spectral analysis, the cumulative squared-singular-value mass
of the feature matrix and the cumulative absolute eigenvalue mass of $W$ are recorded. Both
masses describe how quickly the total is accumulated as components are added.

\subsubsection{Exploratory comparison}
\label{subsubsec:exploratory-comparisons}

\paragraph{\textbf{Cognitive and sex subgroups.}}
\label{par:cognitive-sex-subgroups}
In this exploratory analysis, discrimination is examined across cognitive ability and sex using out-of-fold predictions from the primary cohort. Predictions are first matched to the correct phenotypic records after participant labels and site order are verified. Cognitive effects are then examined in three ways. First, full-scale IQ is centred within site and divided into quartiles, and AUROC is calculated within each quartile so that the groups reflect IQ relative to the site mean. Second, autistic participants are divided into full-scale IQ bands below 85, 85--99, 100--114, and 115 or above, while the same 468 non-autistic participants are used as the comparison group for every band. Third, autistic and non-autistic participants are matched on full-scale IQ within site using greedy nearest-neighbour matching without replacement and a maximum difference of five IQ points. Unmatched autistic participants are excluded, and AUROC is then estimated within site-centred IQ quartiles of the matched sample. These three cognitive-subgroup designs are illustrated in Figure~\ref{fig:designs}d. Additional analyses examine the relationship between classifier scores and full-scale, verbal, and performance IQ separately within autistic and non-autistic participants. Before Spearman correlation is calculated, site, age, and mean framewise displacement are regressed from both the classifier score and IQ measure. Finally, sex-specific AUROCs are calculated only for subgroups containing at least 20 participants and both diagnostic classes, reducing instability from very small or single-class subsets.

\paragraph{\textbf{Cohort definition.}}
\label{par:cohort-definition}
This exploratory sensitivity analysis examines whether model performance changes with cohort definition. The static tangent classifier and the primary three-member ensemble are evaluated first in the primary cohort of 871 participants and then in the wider quality-control cohort of 1,032 participants. The model specification is kept unchanged across both cohorts so that any performance difference reflects the cohort definition rather than a change in modelling choices. The static tangent classifier is evaluated with pooled ten-fold validation, while the primary ensemble is evaluated with both pooled ten-fold and leave-one-site-out validation.

\paragraph{\textbf{Scan duration and sample size.}}
\label{par:duration-sample-size}
The exploratory duration analysis examines how performance changes as the amount of scan data is reduced. Participants with at least 232 volumes are retained, and the same scans are truncated to 40, 60, 90, 120, 160, 200, and 232 volumes. This keeps the participant sample fixed across durations. Covariance matrices are re-estimated at each duration, and the analysis is repeated with additional shrinkage values of 0.15, 0.30, and 0.50 to account for the reduced number of volumes available in shorter scans. A separate sample-size analysis examines performance using full-length scans with participant subsets of 100, 200, 300, 400, 550, 700, and 825. One unstratified random subset is used at each size, and additional shrinkage is fixed at $r=0.15$ so that only sample size is varied. Sample size is also varied within the New York University (NYU) site using five-fold validation, allowing this effect to be examined without between-site variation. Finally, the association between scan duration or sample size and AUROC is summarised using Spearman correlation (Section~\ref{subsubsec:correlation-analyses}).

\paragraph{\textbf{Training-cohort homogeneity.}}
\label{par:training-homogeneity}
In this exploratory comparison, the effect of training-cohort homogeneity on AUROC is
examined. Here, homogeneity refers to whether the training participants come from one site
or from several sites. For this purpose, NYU is used as the single-site source, and a static
tangent model with additional shrinkage $r=0.30$ is fitted. Four conditions are then formed
by crossing one-site or multi-site training with same-site or other-site testing. In each
condition, 120 training participants and 40 test participants are sampled without class
stratification over 20 random repetitions. The sample sizes are held equal across conditions
so that the conditions differ only in the sites that supply the participants. Across all
four conditions, AUROC is used as the outcome.

\paragraph{\textbf{Clinical-instrument concordance.}}
\label{par:instrument-concordance}
This exploratory analysis examines agreement between the ADI-R and ADOS criteria within autistic participants who have complete scores for both instruments. The ADI-R criterion requires scores of at least 10 for social interaction, 8 for verbal communication, and 3 for restricted or repetitive behaviour, while the ADOS criterion requires a total score of at least 7 \citep{Lordetal1994,Lordetal2000}. Agreement is summarised using the four cells of the resulting 2 $\times$ 2 table, the discordance rate, and Cohen's $\kappa$, which measures agreement beyond chance. The ABIDE-I diagnostic label remains the classification outcome throughout. The concordance results are also combined with Social Responsiveness Scale data from non-autistic participants to examine the effect of possible label noise under the mutually contaminated model of \citet{Menonetal2015}. The autistic-label noise rate, $\varepsilon_1$, is initially set to half the observed disagreement rate $\eta$, assuming that each discordant case reflects an error in one of the two instruments. The non-autistic noise rate, $\varepsilon_0$, is estimated from the 137 non-autistic participants with available SRS totals as the proportion whose raw score is at least 60. With $\pi_1$ and $\pi_0$ denoting the proportions of autistic and non-autistic participants in the primary cohort, the corresponding accuracy ceiling is $1-(\pi_0\varepsilon_0+\pi_1\varepsilon_1)$. The observed AUROC and accuracy of the primary ensemble are also adjusted using these noise rates. The corrected AUROC is $(\mathrm{AUROC}-0.5)/(1-\varepsilon_0-\varepsilon_1)+0.5$, while corrected accuracy is $(\mathrm{accuracy}-(\varepsilon_0+\varepsilon_1)/2)/(1-\varepsilon_0-\varepsilon_1)$. The latter expression is exact for balanced accuracy but is only approximate for accuracy because the diagnostic groups are unequal in size. These noise rates are not identifiable from the available data. In particular, two diagnostic instruments applied to a single population do not uniquely determine their sensitivities, specificities, and prevalence, and no participant is negative on both instruments. In addition, the conventional SRS threshold of 60 applies to a T-score, whereas the release provides raw totals. For these reasons, the accuracy ceiling and corrected performance are reported as a range over $\varepsilon_1\in\{0,\eta/2,\eta\}$ and raw-total thresholds of 40, 50, 55, 60, and 70.

\paragraph{\textbf{Coefficient and SHAP summaries.}}
\label{par:coefficient-shap}
In this exploratory analysis, the model coefficients are treated as descriptive parameters.
A coefficient reflects the contribution of a feature to the prediction given all other
features, rather than a group difference in that feature. For this reason, the
interpretation of group differences is based on the residualised analyses in
Section~\ref{subsec:group-differences}.

Shapley additive explanations (SHAP) are then calculated to describe how each feature
contributes to individual predictions. Under SHAP, each feature is assigned a share of the
difference between a participant's model output and the average output over a background
set. For a linear model and a training-fold background, the SHAP values are calculated with
Equation~\eqref{eq:shap-value}:
\begin{equation}
\phi_{ij}
=
w_j
\left(
z_{ij}
-
\mathbb{E}_{\mathrm{T}}[z_j]
\right).
\label{eq:shap-value}
\end{equation}

In Equation~\eqref{eq:shap-value}, $\phi_{ij}$ denotes the SHAP value of feature $j$ for participant $i$, $w_j$ is the corresponding model coefficient, $z_{ij}$ is the standardised feature value, and $\mathbb{E}_{\mathrm{T}}[z_j]$ is the training-fold mean of that feature. SHAP values are computed out of fold using $C=0.01$ under pooled ten-fold validation, with the corresponding training partition used as the explainer background \citep{LundbergLee2017}. This keeps held-out participants separate from the explanation process. Within each fold, the computed SHAP values are also compared with the closed-form values using Pearson correlation to verify agreement. For interpretation, mean absolute SHAP values are first calculated for individual edges so that positive and negative contributions do not cancel. These values are then aggregated into Yeo-7 system pairs using system labels assigned from the parcel centroids. This assignment differs from the voxel-overlap mapping used in Section~\ref{subsec:group-differences}. Unassigned edges are omitted, while within-system edges contribute twice to diagonal cells. Finally, system-level summaries and three participant-level waterfall examples are retained to illustrate how feature contributions combine to produce individual model outputs.

\paragraph{\textbf{Component--symptom associations.}}
\label{par:component-symptom}
In this exploratory analysis, a rank-20 spectral classifier is fitted once to the complete primary cohort so that all participants are represented using the same set of components. The tangent reference, scaler, and classifier are therefore estimated from the full cohort and its diagnostic labels, while symptom scores are used only in the subsequent association analysis and do not influence component fitting. Within autistic participants, each component score is related to the symptom measures after age, sex, full-scale IQ, and site are regressed from both variables. Spearman correlation is then calculated between the resulting residuals, and p-values are adjusted across all component--instrument combinations. Because these associations rely on components obtained from a single full-cohort fit, component stability is also evaluated. Components estimated within each cross-validation training fold are matched to those from the full cohort, and stability is summarised by the highest absolute cosine similarity for each component. The absolute value is used because the sign of an eigenvector is arbitrary.

\subsection{Exploratory group-difference analyses}
\label{subsec:group-differences}

\subsubsection{Connection construction and nuisance adjustment}
\label{subsubsec:connection-adjustment}

The group-difference analyses directly compare functional connectivity between autistic and non-autistic participants. These exploratory analyses use the group-contrast cohort of 863 participants and 196 retained regions. Pairwise Pearson correlations are Fisher transformed and vectorised before analysis. For each connection, a linear nuisance model is fitted using age, sex, mean framewise displacement, and acquisition-site indicators. The resulting adjusted connectivity values are then used for the group comparisons. Missing nuisance values are replaced with the median of the analysis cohort so that all participants \reviewhl{are retained}.


Standardised group differences are then calculated from the residual connection values.
For connection $e$, the effect is defined by Equation~\eqref{eq:group-difference}:
\begin{equation}
d_e
=
\frac{
\overline{R}_{e,\mathrm{autistic}}
-
\overline{R}_{e,\mathrm{comparison}}
}{
\operatorname{SD}(R_e)
}.
\label{eq:group-difference}
\end{equation}

In Equation~\eqref{eq:group-difference}, $\operatorname{SD}$ is the standard deviation of the residuals across both groups, while the barred terms denote the corresponding group means. The comparison group represents non-autistic participants, so negative values indicate lower connectivity in autistic participants. Standardising by the pooled residual variation places all connection effects on a common scale, allowing them to be compared and averaged.


\subsubsection{Exploratory functional-system contrasts}
\label{subsubsec:system-contrasts}

Connection-level effects are aggregated into functional-system contrasts by grouping connections according to the systems of their two endpoints. For each system pair with at least 20 connections, the standardised connection effects are averaged to obtain a system-level effect. Statistical significance is then assessed by permuting diagnostic labels within site. The nuisance model and its residuals are fitted once before permutation, and only the group contrast is recomputed for each reshuffled label set. This keeps the nuisance adjustment fixed while testing whether the observed system-level effect is larger than expected under relabeling.

Because diagnosis is associated with age, sex and head motion, permuting the diagnosis labels would compare the observed statistic with a null in which diagnosis is independent of those covariates, which is not the covariate-adjusted null that the claim requires. The system-level and regional tests are therefore repeated under the Freedman--Lane procedure. A reduced model containing only the nuisance variables is fitted first, and its residuals are permuted within site and added back to the reduced fitted values. The full model, in which the diagnosis vector retains its true position and therefore its association with the covariates, is then refitted to each permuted outcome, and the diagnosis coefficient is recorded. Five thousand permutations are drawn under seed 42, two-sided p-values are formed by comparing the absolute observed statistic with its permutation distribution, and the Benjamini--Hochberg procedure is applied across the system pairs and across the somatomotor regions separately.


\subsubsection{Exploratory motion-matched sensitivity analysis}
\label{subsubsec:motion-matching}

The sensitivity of the group contrasts to head motion is examined using motion-matched samples. Autistic and non-autistic participants are matched within site by mean framewise displacement using greedy matching without replacement. Matching within site keeps acquisition site balanced between the two groups. Before matching, missing nuisance values are imputed using medians from the complete pre-matching cohort so that imputation does not depend on the final matched sample. The group-contrast and motion-matching procedures are illustrated in Figure~\ref{fig:designs}c.

Two matching calipers are used to balance motion similarity against retained sample size. For the system-level analysis, the maximum allowed difference in mean framewise displacement is 0.01 mm, producing 256 matched pairs (512 participants). For the regional analysis, a 0.02 mm caliper is used, producing 280 pairs (560 participants). Nuisance adjustment and site-stratified permutation testing are then repeated within each matched cohort. Thus, the matched analyses use the same group-contrast procedure as the full cohort but with different participant samples. The difference in mean framewise displacement between groups is also tested with the Mann--Whitney U test before and after matching to determine how effectively the matching reduces motion imbalance.

The relation between head motion and overall connectivity is examined separately. For each participant, mean connection strength is calculated as the average Fisher-transformed correlation across the 19,110 retained connections. Its association with mean framewise displacement is then estimated using Spearman correlation across the group-contrast cohort. The analysis is repeated after acquisition site is regressed from both variables. This site-adjusted estimate is included to reduce the influence of site-level differences in both motion and connectivity.

\subsubsection{Exploratory regional analysis and internal stability}
\label{subsubsec:regional-analysis}

In this exploratory analysis, regional effects are examined within a selected functional system, with the somatomotor system used as the default. For each parcel, the regional effect is calculated as the mean standardised group difference across its within-system connections, and p-values are adjusted across the parcels tested within that system. Because this analysis is performed after inspection of the system-level results, it is treated as exploratory. The stability of the resulting regional ranking is then assessed using 200 random split-half repetitions. In each repetition, the cohort is divided into two independent halves, and nuisance adjustment and regional effects are refitted separately in each half. Agreement between the two regional rankings is measured using Spearman correlation. The same procedure is also applied to all within-system connections and to the top 50 connections. As an additional reference, 100 randomly selected region sets of the same size are evaluated using the first 50 split-half repetitions, allowing the stability of the target system to be compared with that expected from similarly sized random sets.

\subsection{Statistical analysis}
\label{subsec:statistical-analysis}

\subsubsection{Analysis status and multiplicity}
\label{subsubsec:analysis-status}

All prediction and group-difference analyses are interpreted within a common evidential framework. No confirmatory hypothesis-testing family is predefined. Instead, analyses are classified as primary descriptive, secondary, or exploratory, and this status is stated in each analysis heading. The primary descriptive results include the primary-cohort ensemble, the matched-shrinkage comparison of static and windowed connectivity, the feature-selection scope analysis, and the pooled and within-participant symptom associations. Pairwise model comparisons, including those within the primary descriptive analyses, are therefore interpreted descriptively rather than as confirmatory hypothesis tests.

Multiplicity is handled according to the type of analysis. Pairwise comparison p-values are reported without adjustment and are presented together with effect sizes. Benjamini--Hochberg q-values are also provided across these pairwise comparisons for reference, but they are not used as a decision rule. In contrast, Benjamini--Hochberg false-discovery-rate correction is applied to families of related network, regional, symptom-association, and component--symptom tests. For the pooled and within-participant symptom analyses, the six symptom measures form one testing family within each combination of validation protocol and correlation scope.

\subsubsection{Paired predictive comparisons}
\label{subsubsec:paired-comparisons}

Conditions evaluated on the same participants are compared using their stored out-of-fold predictions. Differences in AUROC are tested with DeLong's test, while differences in binary decisions at the 0.5 threshold are assessed with McNemar's test. These paired tests are used to compare discrimination and threshold-based classification under the same participant sample.

Interval estimates are additionally obtained from a site-clustered bootstrap, because DeLong
and Wilson intervals computed on pooled out-of-fold predictions condition on one fitted split
structure and treat participants from the same acquisition site as independent. In this
procedure, the 20 site labels are resampled with replacement 2,000 times using random seed 42,
all participants of a drawn site are retained together, and the performance measure is
recomputed within each resample. Percentile intervals are then taken across resamples. Because
whole sites rather than individual participants are resampled, the resulting intervals carry
the dependence induced by site clustering. These intervals are reported alongside the DeLong
and Wilson intervals, and the parametric intervals are interpreted as lower bounds on
uncertainty rather than as complete accounts of it.


\subsubsection{Permutation analyses}
\label{subsubsec:permutation-analyses}

In the functional-system and regional analyses, 5,000 permutations of diagnostic labels are
performed within site. Labels are permuted within site so that each permuted dataset keeps
the number of autistic and non-autistic participants at every site. The two-sided p-value
includes a plus-one adjustment and is given by Equation~\eqref{eq:permutation-p}. Here
$T_{\mathrm{obs}}$ is the observed system- or region-level effect, and $T_b^\ast$ is the
corresponding effect after permutation $b$:
\begin{equation}
p
=
\frac{
1
+
\sum_{b=1}^{B}
\mathbb{1}\left(
\left|T_b^\ast\right|
\geq
\left|T_{\mathrm{obs}}\right|
\right)
}{
B+1
},
\qquad B=5000.
\label{eq:permutation-p}
\end{equation}

In Equation~\eqref{eq:permutation-p}, $B$ is the number of permutations, and $\mathbb{1}(\cdot)$
equals one when its condition holds and zero otherwise. The equation adds one to the
numerator and denominator, because the observed labelling is counted as one of the
possible labellings. As a result, the smallest attainable value at $B=5000$ is $1/5001$,
and values at that floor are reported as $p < 2\times10^{-4}$.

\subsubsection{Correlation analyses}
\label{subsubsec:correlation-analyses}

Spearman rank correlation is used in the symptom, cognitive-subgroup,
motion--connectivity, internal-stability, and scaling analyses. This coefficient is used
because it measures monotonic association without assuming a linear relation or normally
distributed variables. In addition, each coefficient is accompanied by its participant
scope, available sample size, and adjustment procedure. These details are reported because
the participant set and the adjustment differ between analyses.

\subsection{Reproducibility and scope}
\label{subsec:reproducibility-scope}

All analyses are implemented with fixed random seeds and stored outputs so that the main experiments, participant splits, and prediction results \reviewhl{are reproducible}. Fold-level results and out-of-fold predictions are retained, and participant labels and site information are checked when predictions are linked with phenotypic variables or used in paired comparisons. The scope of the findings is defined by the datasets and analysis settings used in each experiment. Primary diagnostic-group discrimination is evaluated in the 871-participant primary ABIDE-I cohort, while the wider quality-control cohort of 1,032 participants is used for cohort sensitivity analyses. Group-difference analyses use 863 participants and 196 retained regions. Generalisation is assessed across the acquisition sites represented in ABIDE-I, and the results therefore remain specific to this multisite research setting. Dynamic analyses are limited to the 50-volume, six-window design, and their temporal duration varies across sites because repetition time is not harmonised. Group-level findings are also conditional on the selected preprocessing pipeline, CC200 parcellation, nuisance adjustment, and recorded mean framewise displacement. These design choices define the conditions under which the reported results \reviewhl{are interpreted}.

The computational environment is recorded with the results. The analyses reported here are produced under Python 3.12.8 with NumPy 2.4.6, pandas 3.0.3, SciPy 1.18.0 and scikit-learn 1.8.0 on macOS (arm64), using random seed 42 throughout. The imaging input is the C-PAC \texttt{filt\_noglobal} derivative of the ABIDE Preprocessed Initiative, and the parcellation file is the CC200 atlas volume distributed with that release. The 871 participant identifiers of the primary cohort are stored with the results and are summarised by the SHA-256 digest \texttt{6c943ff4\ldots9b81a057} of their sorted list, so that the cohort can be checked without redistributing identifiers. Fold assignments are reproducible rather than stored as free text: the pooled folds are the test partitions of a stratified ten-fold split with shuffling under seed 42, and the leave-one-site-out order follows the site labels in alphabetical order. Both orders were verified against the stored labels and site codes before the out-of-fold predictions were linked to phenotypic variables. Code and result files are prepared for deposit in a public repository under the ABIDE data-use terms.
